\chapter{PENELITIAN TERKAIT}

\section{Kubernetes}

Kubernetes adalah platform orkestrasi kontainer \textit{open-source} yang dirancang untuk mengotomatisasi deployment, penskalaan, dan pengelolaan aplikasi kontainer \citep{kubernetes2022}. Kubernetes menyediakan abstraksi yang kuat untuk mengelola kompleksitas infrastruktur modern dengan konsep seperti Pod, Deployment, Service, dan Namespace.

\subsection{Arsitektur Kubernetes}

Arsitektur Kubernetes terdiri dari:

\begin{itemize}
    \item \textbf{Control Plane:} Mengelola status klaster, termasuk API Server, etcd, Scheduler, dan Controller Manager.
    \item \textbf{Worker Nodes:} Menjalankan aplikasi kontainer menggunakan kubelet dan kube-proxy.
\end{itemize}

\subsection{Manifest Kubernetes}

Semua objek Kubernetes didefinisikan secara deklaratif menggunakan manifest YAML. Pendekatan deklaratif ini menjadi fondasi bagi praktik GitOps, di mana status yang diinginkan dari sistem dideskripsikan dalam file konfigurasi yang tersimpan di repositori Git \citep{kubernetes-docs}.

\section{GitOps}

GitOps adalah paradigma operasional yang menggunakan Git sebagai sumber kebenaran tunggal untuk infrastruktur dan aplikasi. Konsep ini pertama kali dipopulerkan oleh Weaveworks pada tahun 2017 \citep{weaveworks-gitops}.

\subsection{Prinsip GitOps}

Menurut OpenGitOps Working Group, terdapat empat prinsip utama GitOps \citep{gitops-principles}:

\begin{enumerate}
    \item \textbf{Declarative:} Sistem dideskripsikan secara deklaratif.
    \item \textbf{Versioned and Immutable:} Status yang diinginkan disimpan dalam versi yang tidak dapat diubah.
    \item \textbf{Pulled Automatically:} Software agent secara otomatis menarik perubahan dari sumber kebenaran.
    \item \textbf{Continuously Reconciled:} Software agent terus-menerus memantau dan mengoreksi penyimpangan.
\end{enumerate}

\subsection{Manfaat GitOps}

Implementasi GitOps memberikan berbagai manfaat, antara lain:

\begin{itemize}
    \item Audit trail lengkap untuk setiap perubahan sistem.
    \item Kemudahan \textit{rollback} ke versi sebelumnya.
    \item Kolaborasi yang lebih baik antara tim dengan menggunakan alur kerja \textit{pull request}.
    \item Peningkatan keamanan melalui kontrol akses berbasis Git.
\end{itemize}

\section{ArgoCD}

ArgoCD adalah \textit{GitOps controller} \textit{open-source} untuk Kubernetes yang mengotomatiskan deployment aplikasi yang diinginkan ke target environment \citep{argocd-docs}. ArgoCD mendukung berbagai format konfigurasi, termasuk Kustomize, Helm, dan plain YAML.

\subsection{Fitur Utama ArgoCD}

Beberapa fitur utama ArgoCD antara lain:

\begin{itemize}
    \item \textbf{Automated Sync:} Sinkronisasi otomatis dari Git ke klaster.
    \item \textbf{Sync Waves:} Mengatur urutan deployment berdasarkan dependensi.
    \item \textbf{Resource Health:} Monitoring status kesehatan aplikasi.
    \item \textbf{Rollback:} Kemudahan mengembalikan ke versi sebelumnya.
\end{itemize}

\section{Tinjauan Keamanan}

Keamanan dalam implementasi GitOps menjadi perhatian penting \citep{gitops-security}. Beberapa aspek keamanan yang relevan dengan penelitian ini meliputi:

\begin{itemize}
    \item \textbf{Secrets Management:} Penggunaan Sealed Secrets untuk menyimpan rahasia terenkripsi di Git.
    \item \textbf{Network Policies:} Kontrol lalu lintas jaringan antara Pod dalam klaster.
    \item \textbf{Pod Security:} Konfigurasi keamanan kontainer seperti non-root user dan \textit{read-only root filesystem}.
\end{itemize}

\section{InvenioRDM}

InvenioRDM adalah platform manajemen data riset \textit{open-source} yang dibangun di atas framework Invenio. Platform ini menyediakan solusi lengkap untuk mengelola, mempublikasikan, dan memelihara data riset dengan fitur pencarian, depositori, dan manajemen akses.
